% Source: sources/returns/S6_04_MAPS_MANIFOLDS/sections/manifolds.tex:4717–4842
\paragraph{Theorem D.2.23 (Proved: uniform all-order fundamental groups on the natural catalogue).}
For every one of the $172$ natural initial rows of D.2.22, retain any
of the stated original or Mumford fillings and perform any admissible
free affine logarithmic replacements. At an odd $I_n^*$, a logarithmic
replacement here means the free cyclic quotient of its rank-one
semistable Mumford cover. Regular logarithmic sites are allowed.
Let $m_i\ge1$ be their orders, $v_i$ their integral power vectors in
$H_1(T^4;\mathbb Z)$, and put $n_i=\gamma^*(v_i)$.
For a chosen system of peripheral lifts, let $e$ be the integer in
the global relation of their surviving $\gamma$ coordinates.
Then the complete fundamental group is
\[
 \pi_1(X)=\left\langle \gamma,c_1,\ldots,c_s\ \left|\
 [\gamma,c_i]=1,\ c_i^{m_i}=\gamma^{n_i},\
 c_1\cdots c_s=\gamma^e\right.\right\rangle .                 \tag{D.FG1}
\]
The canonical lifts from the construction, obtained by unrolling the
translation quotient before filling, have $e=0$. Changing $c_i$ to
$c_i\gamma^{t_i}$ replaces $(n_i,e)$ by
$(n_i+m_it_i,e+\sum t_i)$ and leaves (D.FG1) unchanged.
Thus the formula also applies directly to other marked twist conventions.
The first integral homology is the Smith cokernel of the $(s+1)$ square
presentation matrix with rows
\[
 (-n_i,0,\ldots,m_i,\ldots,0)\quad(1\le i\le s),
 \qquad (-e,1,\ldots,1).                                    \tag{D.FG2}
\]
There is no multiplicity bound and no restriction to invariant rational
representatives of the affine translations: their actual power vectors
are used.

After eliminating order-one sites by $c_i=\gamma^{n_i}$, replace $e$
by $e-\sum_{m_i=1}n_i$. If at least three sites of order greater than
one remain, the group is nontrivial. With zero, one, or two remaining
sites, it is abelian and is trivial precisely when, respectively,
\[
 |e|=1,\qquad |m_1e-n_1|=1,\qquad
 |m_2n_1+m_1n_2-em_1m_2|=1.                                \tag{D.FG3}
\]
In the two-site case its invariant factors are the gcd of the entries
of the two-by-two relation matrix and its determinant divided by that
gcd; when the determinant vanishes use one free summand. In particular,
simple connectivity forces coprime remaining multiplicities and
$\gcd(m_i,n_i)=1$ at every remaining site.

\begin{proof}
We supply the integral vanishing comparison, since rational monodromy
alone would not prove this statement. The exact program
\texttt{manifolds\_sphere\_gluing.py} constructs all $172$ integral
monodromy tuples. Every zero-free point is locally narrow for at least
one of $P,Q$, so its lift is the one already proved in D.2.19.
The matrix at the one component-supported zero is independently fixed
by the global product relation. It agrees integrally with the normalized
weight $1/n$ derived in D.2.22. The program verifies every product,
all local Smith forms, and the preserved alternating form.

At a multiplicative Mumford cap the primitive vanishing sublattice is
$\operatorname{sat}\operatorname{im}(A_i-I)$: its individual primitive
circles bound the toric vanishing thimbles. The zero cap has rank two.
Every other multiplicative cap in all $172$ rows has rank one and
Smith width one; there is no nonprimitive necklace hidden in this
assertion. If instead the original zero-free quotient is retained,
these same circles are still nullhomotopic. When the line multidegree
is zero, the local line trivialization identifies the elliptic vanishing
thimble with this primitive circle. When it is nonzero, the natural
line has component degree vector $e_{P\text{-component}}-e_0$;
its entries have gcd one, so the component two-spheres kill the
$\delta$ circle integrally. The elliptic vanishing thimble then kills
its primitive elliptic coordinate as well, and hence kills the required
mixed circle. These comparisons do not replace an integral image by
its saturation without a geometric bounding disc.

Close these vanishing lattices under all monodromy matrices. Exact
unimodular operations show that their normal closure is
$\ker\gamma^*$ in $170$ rows. In the two rows
\[
 (46/II,a,b)=(46/II,1,4),\ (46/II,4,1)
\]
its quotient in $\ker\gamma^*$ is $\mathbb Z/3$. The $I_3^*$ cap
kills this last quotient. An original cap kills its elliptic vanishing
loops. For a semistable logarithmic cap, take its cyclic base cover.
The upstairs rank-one Mumford model has an actual disc bounding the
primitive circle in $\operatorname{im}(A_*^2-I)\otimes\mathbb Q$.
A selected regular fibre upstairs maps diffeomorphically to the
regular fibre downstairs: all nonidentity deck transformations move
its base root. Projecting this disc therefore bounds the same primitive
fibre circle downstairs, rather than a multiple. Appending this column
gives three unit Smith factors in each exceptional row. This proves
that all the fibre subgroup except the primitive $\gamma$ is killed,
for every allowed filling choice.

Over the punctured base the torus bundle has a section up to homotopy,
since the base retracts onto a graph. Its group is the corresponding
semidirect product. Modulo the vanishing normal closure just proved,
every monodromy acts trivially and its fibre group is exactly
$\mathbb Z\gamma$. A free affine log cap has a cyclic cover whose
fundamental group is its torus lattice, modulo the primitive vanishing
line in the semistable case. The covering extension has presentation
$c_i\lambda c_i^{-1}=A_i\lambda$, $c_i^{m_i}=v_i$; no further relation
is introduced. After the preceding quotient this becomes exactly
$[c_i,\gamma]=1$, $c_i^{m_i}=\gamma^{n_i}$.

For an original cap the translation quotient can be unrolled, and the
remaining group is $\mathbb Z\gamma$, with the base meridian trivial.
The same unrolling extends across a Mumford cap by unrolling its long
period in the infinite toric model. These are restrictions of the
original unrolled translation construction, so their global peripheral
relation has zero $\gamma$ coordinate. Equivalently, transition
functions of the line bundle contribute only the $\delta$ coordinate.
This proves $e=0$ in the canonical marking. Other meridian choices
produce exactly the stated change of $(n_i,e)$. Van Kampen now gives
(D.FG1), including its completeness, and abelianization gives (D.FG2).

Quotienting (D.FG1) by its central fibre yields the sphere-orbifold
presentation with cone orders $m_i>1$. With at least three such orders
it is nontrivial. One direct verification is the polygon-reflection
construction: a polygon with angles $\pi/m_i$ exists in the sphere,
plane or hyperbolic plane according to the angle sum, and products of adjacent side reflections give rotations of exact
orders $m_i$ whose cyclic product is the identity. They therefore
give a nontrivial representation of the orbifold presentation. The familiar finite spherical triangle
cases have the same nontrivial rotations. For at most two sites,
eliminate the product relation; the remaining generators commute.
The zero-, one-, and two-site integral relation matrices give (D.FG3).
Reducing the last determinant modulo each multiplicity proves the
final coprimality assertions.
\end{proof}
